% !PRINTED 430
\weberpage{430}

In the chain of compositions
\weq{3}{P,\ P^2,\ P^3,\ \ldots}
there may occur a form that belongs to the \gesperrt{principal class}, that is, one equivalent to $(1,\,0,\,\tfrac14 m)$ or $[1,\,1,\,\tfrac14(m+1)]$. If this happens for $P^{\varepsilon}$, then $p^{\varepsilon}$ is representable by the principal form, and there is a \gesperrt{proper} representation
\weq{4}{4\,p^{\varepsilon} = y^2 - D\,x^2.}
If conversely such a representation is possible, then there is a form $(p^{\varepsilon},\,b,\,c)$ in the principal class.

\gesperrt{If $\varepsilon$ is the least positive exponent for which equation (4) is properly solvable, then $\varepsilon$ is called the index of $p$, and $p$ belongs to the exponent $\varepsilon$.}

If
\weq{5}{P = (p,\,b,\,p^{\varepsilon-1}c),}
then
\[
b^2 \equiv D \pmod{4\,p^{\varepsilon}},
\]
and when this condition is satisfied, then in the form $(p,\,b,\,c')$ of discriminant $D$ the third coefficient is divisible by $p^{\varepsilon-1}$. The form (5) is preserved if we replace $b$ by any number $b'$ congruent to it with respect to the modulus $2\,p^{\varepsilon}$. If therefore $a$ is not divisible by $p$, and
\[
\varphi = (a,\,B,\,C)
\]
is any form of discriminant $D$ concordant with $P$, then the two congruences
\[
x \equiv b \pmod{2\,p^{\varepsilon}}, \quad x \equiv B \pmod{2\,a}
\]
can be satisfied at the same time, and if one then puts this $x$ in place of $b$ and $B$, it follows that:

III. \gesperrt{If $l$ is the class of $P$ and $k$ an arbitrary class of discriminant $D$, then in $l$ and $k$ one can choose the representatives:}
\weq{6}{(p,\,b,\,a\,c\,p^{\varepsilon-1}),\quad (a,\,b,\,c\,p^{\varepsilon}),}
and in the compounded class $lk$ one obtains a representative
\weq{7}{(a\,p,\,b,\,p^{\varepsilon-1}c).}
If therefore
\[
\omega = \frac{-b+\sqrt{-m}}{2\,a}
\]
is a root of the class $k$, then $\omega\colon p$ is a root of the class $lk$.
